\appendix

\stepcounter{appendixnum}
\chapter*{Приложение}
\addcontentsline{toc}{chapter}{Приложение}
\setcounter{section}{0}
\setcounter{subsection}{0}
\setcounter{subsubsection}{0}
\setcounter{table}{0}
\renewcommand{\thesection}{\Alph{section}}
\renewcommand{\thesubsection}{\thesection.\arabic{subsection}}
\renewcommand{\thesubsubsection}{\thesubsection.\arabic{subsubsection}}
\renewcommand{\thetable}{\thesection.\arabic{table}}

\section{Примеры использования системы}

\subsection{Установка и запуск}

\subsubsection{Требования}

\begin{itemize}
\item Python 3.10+
\item CUDA-совместимый GPU (рекомендуется 6GB+ VRAM)
\item 8GB+ RAM
\item 5GB свободного места на диске
\end{itemize}

\subsubsection{Установка зависимостей}

\begin{verbatim}
pip install torch torchvision torchaudio --index-url
  https://download.pytorch.org/whl/cu118
pip install -r requirements.txt
\end{verbatim}

\subsubsection{Подготовка данных}

\begin{verbatim}
# Скачать POP909 датасет
python scripts/download_pop909.py

# Подготовить FASTA файлы
python scripts/prepare_fasta.py
\end{verbatim}

\subsection{Обучение модели}

\subsubsection{Конфигурация}

Пример конфигурационного файла \texttt{config.json}:

\begin{verbatim}
{
  "model": {
    "d_model": 384,
    "n_layers": 6,
    "n_heads": 6,
    "dim_feedforward": 1536,
    "dropout": 0.20
  },
  "training": {
    "batch_size": 8,
    "gradient_accumulation_steps": 4,
    "learning_rate": 0.0003,
    "weight_decay": 0.05,
    "max_epochs": 35
  }
}
\end{verbatim}

\subsubsection{Запуск обучения}

\begin{verbatim}
python train_bio_music_v2.py \
  --config configs/pipeline_v2_medium_rtx2060_fast.json \
  --output_dir results/my_experiment
\end{verbatim}

\subsection{Генерация музыки}

\subsubsection{Из FASTA файла}

\begin{verbatim}
python generate_from_fasta_v2_fragmented.py \
  --fasta data/fasta/training/thesis_reference_genomes.fa \
  --checkpoint results/v2_medium_rtx2060_fast/
    checkpoints/structured_pipeline.pt \
  --config configs/pipeline_v2_medium_rtx2060_fast.json \
  --bars-per-fragment 4 \
  --output generated_music.mid
\end{verbatim}

\subsubsection{Через веб-интерфейс}

\begin{verbatim}
cd web
python app.py
\end{verbatim}

Откройте браузер по адресу \texttt{http://localhost:5001}.

\subsection{Примеры входных данных}

\subsubsection{FASTA формат}

\begin{verbatim}
>NC_000913.3 Escherichia coli K-12 MG1655
AGCTTTTCATTCTGACTGCAACGGGCAATATGTCTCTGTGTGGATTAAAAAAAGAGTGTC
TGATAGCAGCTTCTGAACTGGTTACCTGCCGTGAGTAAATTAAAATTTTATTGACTTAGG
TCACTAAATACTTTAACCAATATAGGCATAGCGCACAGACAGATAAAAATTACAGAGTAC
\end{verbatim}

\subsubsection{Protein FASTA}

\begin{verbatim}
>sp|P0A7B8|RPOB_ECOLI DNA-directed RNA polymerase
MQTLNDRLLVLPVELLRGETVRRTQTDSPYAQLADALRQRLAERLGVQPLVVFEMADLR
AAFNHSSQPGTTLPQYVEMLPKQKVLGAMVTHRLSAEEKAELWQRIRQNVREARKQIVR
\end{verbatim}

\section{Структура проекта}

\subsection{Основные модули}

\begin{verbatim}
biosonification/
├── bio_music_pipeline/
│   ├── v2/
│   │   ├── bio.py              # Биологический энкодер
│   │   ├── structured_music.py # Музыкальное представление
│   │   ├── structured_pairing.py # Pairing модуль
│   │   └── structured_model.py # Нейросетевая модель
│   └── utils/
├── configs/                    # Конфигурационные файлы
├── data/                       # Данные (FASTA, MIDI)
├── results/                    # Результаты экспериментов
├── web/                        # Веб-интерфейс
├── train_bio_music_v2.py      # CLI обучения
└── generate_from_fasta_v2_fragmented.py # CLI генерации
\end{verbatim}

\subsection{Ключевые файлы}

\begin{itemize}
\item \texttt{bio.py} — извлечение биологических признаков
\item \texttt{structured\_music.py} — структуры HarmonyBar и MelodyEvent
\item \texttt{structured\_pairing.py} — сопоставление bio и music сегментов
\item \texttt{structured\_model.py} — Bio-Conditioned Transformer
\item \texttt{train\_bio\_music\_v2.py} — основной скрипт обучения
\item \texttt{generate\_from\_fasta\_v2\_fragmented.py} — генерация MIDI
\end{itemize}

\section{Дополнительные таблицы и графики}

\subsection{Детальная статистика обучения}

\begin{table}[htbp]
\centering
\caption{Динамика обучения harmony model по эпохам}
\begin{tabular}{|c|c|c|}
\hline
\textbf{Epoch} & \textbf{Train Loss} & \textbf{Val Loss} \\
\hline
1 & 0.171 & 0.159 \\
2 & 0.126 & 0.145 \\
3 & 0.109 & 0.147 \\
4 & 0.097 & 0.146 \\
5 & 0.087 & 0.150 \\
6 & 0.080 & 0.156 \\
7 & 0.074 & 0.152 \\
8 & 0.070 & 0.156 \\
9 & 0.065 & 0.154 \\
10 & 0.061 & 0.156 \\
\hline
\end{tabular}
\end{table}

\begin{table}[htbp]
\centering
\caption{Динамика обучения melody model по эпохам}
\begin{tabular}{|c|c|c|}
\hline
\textbf{Epoch} & \textbf{Train Loss} & \textbf{Val Loss} \\
\hline
1 & 0.255 & 0.202 \\
2 & 0.167 & 0.179 \\
3 & 0.131 & 0.171 \\
4 & 0.108 & 0.164 \\
5 & 0.091 & 0.161 \\
6 & 0.079 & 0.158 \\
7 & 0.070 & 0.157 \\
8 & 0.062 & 0.163 \\
9 & 0.056 & 0.157 \\
10 & 0.051 & 0.162 \\
11 & 0.047 & 0.169 \\
12 & 0.043 & 0.162 \\
13 & 0.040 & 0.162 \\
14 & 0.038 & 0.162 \\
15 & 0.036 & 0.169 \\
\hline
\end{tabular}
\end{table}

\subsection{Детальная статистика генерации}

\begin{table}[htbp]
\centering
\caption{Метрики для каждого тестового фрагмента}
\begin{tabular}{|c|c|c|c|c|}
\hline
\textbf{Fragment} & \textbf{Notes} & \textbf{Pitch range} & \textbf{Unique pitches} & \textbf{Chord-tone ratio} \\
\hline
1 & 20 & 17 & 9 & 0.650 \\
2 & 24 & 18 & 8 & 0.708 \\
3 & 16 & 12 & 6 & 0.750 \\
4 & 25 & 14 & 8 & 0.800 \\
5 & 23 & 17 & 9 & 0.783 \\
6 & 20 & 10 & 6 & 0.800 \\
7 & 16 & 10 & 6 & 0.500 \\
8 & 26 & 14 & 11 & 0.538 \\
9 & 21 & 12 & 8 & 0.667 \\
10 & 21 & 16 & 10 & 0.762 \\
11 & 20 & 10 & 8 & 0.650 \\
12 & 18 & 21 & 11 & 0.444 \\
\hline
\textbf{Mean} & \textbf{20.83} & \textbf{14.25} & \textbf{8.33} & \textbf{0.6710} \\
\hline
\end{tabular}
\end{table}

\subsection{Примеры биологических признаков}

\begin{table}[htbp]
\centering
\caption{Биологические признаки для фрагмента E. coli}
\begin{tabular}{|l|c|}
\hline
\textbf{Признак} & \textbf{Значение} \\
\hline
GC-content & 0.508 \\
Entropy & 1.982 \\
Aromaticity & 0.089 \\
Instability index & 42.3 \\
Isoelectric point & 6.8 \\
GRAVY & -0.234 \\
Molecular weight & 65432.1 Da \\
\hline
\end{tabular}
\end{table}

\begin{table}[htbp]
\centering
\caption{Control profile для фрагмента E. coli}
\begin{tabular}{|l|c|}
\hline
\textbf{Параметр} & \textbf{Значение} \\
\hline
Tempo & 0.62 \\
Density & 0.54 \\
Harmony change & 0.71 \\
Register & 0.48 \\
Complexity & 0.66 \\
Mode & 0.52 \\
\hline
\end{tabular}
\end{table}
